\pageoriginale

\clearpage
\SGAExposeBookmark{XX}{El teorema de Griffiths}
\begin{center}
{\large SGA 7 --- EXPOSÉ XX}\par
\vspace{2em}
{\Large\underline{El teorema de Griffiths}}\par
\vspace{1.5em}
{\large por N. Katz}
\end{center}

\section*{0. Introducción}

Demostramos el teorema de Griffiths [1], que proporciona una manera
sistemática (aunque no constructiva) de «encontrar» ciclos algebraicos
que son $\mathbb Q_\ell$-homólogos a cero, pero de los cuales ningún
múltiplo no nulo es algebraicamente equivalente a cero. La demostración
que se da aquí (debida a Grothendieck) es la traducción a términos
puramente algebraicos de la demostración original, más o menos
trascendente, de Griffiths. En ambas, el teorema de irreducibilidad
XVIII~6.7 es un punto clave (cf.~3.1.6 y 3.2).

\section*{1. El formalismo de las clases primitivas}

En todo lo que sigue fijamos un número primo $\ell$. Todos los esquemas
considerados se suponen de característica distinta de $\ell$.
Consideremos la situación

\begin{equation}\tag{1.0.1}
\begin{tikzcd}[column sep=huge,row sep=huge]
X \arrow[rr,"\pi"] \arrow[dr,"f"'] & {} & Y \arrow[dl,"g"] \\
{} & S & {}
\end{tikzcd}
\end{equation}
en la que suponemos

\pageoriginale

\begin{equation}\tag{1.0.3}
f,\ g\quad\text{lisos de dimensiones relativas}\quad d_1\ \text{y}\ d_2,
\end{equation}
y
\begin{equation}\tag{1.0.4}
\pi\quad\text{propio}.
\end{equation}
Como $X$ e $Y$ son lisos sobre $S$, la propiedad de ser propio de $\pi$
nos da un \underline{morfismo de Gysin} (SGA~5, Exposé~IV; la referencia
ulterior queda en blanco en la impresión)
\begin{equation}\tag{1.0.5}
\pi_*:H^q(X,\mathbb Q_\ell(i))\longrightarrow
H^{q-2(d_1-d_2)}(Y,\mathbb Q_\ell(i-d_1+d_2)).
\end{equation}
Localizando sobre $S$ para la topología étale, obtenemos asimismo un
morfismo de haces sobre $S$
\begin{equation}\tag{1.0.6}
\pi_*:R^qf_*\mathbb Q_\ell(i)\longrightarrow
R^{q-2\sisi{(d_1-d_2'}{(d_1-d_2)}}g_*
 \mathbb Q_\ell\bigl(\sisi{i-d_1-d_2}{i-d_1+d_2}\bigr),
\end{equation}
y, aplicando el funtor $H^p(S,-)$, un morfismo
\begin{equation}\tag{1.0.7}
\pi_*:H^p(S,R^qf_*\mathbb Q_\ell(i))\longrightarrow
H^p(S,R^{q-2(d_1-d_2)}g_*\mathbb Q_\ell(i-d_1+d_2)).
\end{equation}

\subsubsection*{Definición 1.0.8.}
Sean $(E_r^{p,q},d_r^{p,q})$ y
$({}'E_r^{p,q},{}'d_r^{p,q})$ dos sucesiones espectrales en una categoría
abeliana dada, y sea $(a,b)\in\mathbb Z\times\mathbb Z$. El desplazamiento
de $({}'E_r^{p,q},{}'d_r^{p,q})$ por $(a,b)$ es la sucesión espectral
$({}''E_r^{p,q},{}''d_r^{p,q})$ dada por
\begin{equation}\tag{1.0.9}
\left\{
\begin{aligned}
{}''E_r^{p,q}&={}'E_r^{p+a,q+b},\\
{}''d_r^{p,q}&=\sisi{d_r^{p+a,q+b}}{{}'d_r^{p+a,q+b}}.
\end{aligned}
\right.
\end{equation}

\pageoriginale

\subsubsection*{1.0.10.}
Un morfismo de sucesiones espectrales de
$(E_r^{p,q},d_r^{p,q})$ en $({}''E_r^{p,q},{}''d_r^{p,q})$ se llama
morfismo de bidegrado $(a,b)$ de $(E_r^{p,q},d_r^{p,q})$ en
$({}'E_r^{p,q},{}'d_r^{p,q})$.

El lema siguiente es una consecuencia inmediata de la teoría del
\underline{morfismo de Gysin} (SGA~5, Exposé~IV).

\subsubsection*{Lema 1.1.}
Existe un morfismo canónico $\pi_*$ de bidegrado
$\bigl(\sisi{a}{0},-2(d_1-d_2)\bigr)$ entre las sucesiones espectrales de
Leray
\begin{equation}\tag{1.1.0}
R^p\alpha_*R^qf_*\mathbb Q_\ell(i)\Longrightarrow
R^{p+q}(\alpha f)_*\mathbb Q_\ell(i)
\end{equation}
y
\begin{equation}\tag{1.1.1}
R^p\alpha_*R^qg_*\mathbb Q_\ell(i-d_1+d_2)\Longrightarrow
R^{p+q}(\alpha g)_*\mathbb Q_\ell(i-d_1+d_2),
\end{equation}
que sobre el término $E_2$ viene dado por (1.0.7), y sobre el límite por
(1.0.5).

\subsubsection*{1.2.}
Consideremos ahora la situación
\begin{equation}\tag{1.2.1}
\begin{tikzcd}[column sep=large,row sep=large]
{} & X\times_S Y \arrow[dl,"\mathrm{pr}_1"']
                  \arrow[dr,"\mathrm{pr}_2"] & {} \\
X \arrow[dr,"f"'] & {} & Y \arrow[dl,"g"] \\
{} & S \arrow[d,"\alpha"] & {} \\
{} & T & {}
\end{tikzcd}
\end{equation}

\pageoriginale

en la que $S,X,f,Y,g$ son como en
$\sisi{(1.0.2)}{(1.0.1)}$--(1.0.4), suponiendo ahora además que $f$ es
propio. Sea
\begin{equation}\tag{1.2.2}
z\in H^a(X\times_S Y,\mathbb Q_\ell(b)).
\end{equation}
Entonces $z$ induce una aplicación, aún denotada por $z$,
\begin{equation}\tag{1.2.3}
z:H^q(X,\mathbb Q_\ell(i))\longrightarrow
H^{q+a-2d_1}(Y,\mathbb Q_\ell(i+b-d_1))
\end{equation}
por la fórmula
\begin{equation}\tag{1.2.4}
z(x)=\mathrm{pr}_{2*}\bigl(z\wedge\mathrm{pr}_1^*(x)\bigr),
\end{equation}
donde $\mathrm{pr}_{2*}$ es el homomorfismo de Gysin. Por localización
étale sobre $S$, $z$ define asimismo
\begin{equation}\tag{1.2.5}
z:R^qf_*\mathbb Q_\ell(i)\longrightarrow
R^{q+a-2d_1}g_*\mathbb Q_\ell(i+b-d_1),
\end{equation}
y, aplicando $H^p(S,-)$, aplicaciones
\begin{equation}\tag{1.2.6}
H^p(S,R^qf_*\mathbb Q_\ell(i))\longrightarrow
H^p(S,R^{q+a-2d_1}g_*\mathbb Q_\ell(i+b-d_1)).
\end{equation}

\subsubsection*{Lema 1.3.}
La clase $z$ define un morfismo de bidegrado
$\bigl(0,\sisi{b-2d_1}{a-2d_1}\bigr)$ entre las sucesiones espectrales de
Leray
\begin{equation}\tag{1.2.7}
H^p(S,R^qf_*\mathbb Q_\ell(i))\Longrightarrow
H^{p+q}(X,\mathbb Q_\ell(i))
\end{equation}
y
\begin{equation}\tag{1.2.8}
H^p(S,R^qg_*\mathbb Q_\ell(i+b-d_1))\Longrightarrow
H^{p+q}(Y,\mathbb Q_\ell(i+b-d_1)),
\end{equation}

\pageoriginale

que sobre los términos $E_2$ viene dado por (1.2.6), y sobre los límites
por (1.2.3).

\noindent\underline{Demostración.}
En vista de (1.2.4), y puesto que la «multiplicación por $z$» define un
morfismo de bidegrado $(0,a)$ entre las sucesiones espectrales de Leray
\begin{equation}\tag{1.2.9}
H^p\bigl(S,R^q\sisi{(\mathrm{pr}_1f)_*}{(f\circ\mathrm{pr}_1)_*}
 \mathbb Q_\ell(i)\bigr)
\Longrightarrow H^{p+q}(X\times_S Y,\mathbb Q_\ell(i))
\end{equation}
\begin{center}
\begin{tikzcd}
{} \arrow[d,"z"'] \\ {}
\end{tikzcd}
\end{center}
\begin{equation}\tag{1.2.10}
H^p\bigl(S,R^{q+a}\sisi{(\mathrm{pr}_1f)_*}{(f\circ\mathrm{pr}_1)_*}
 \mathbb Q_\ell(i+b)\bigr)
\Longrightarrow H^{p+q+a}(X\times_S Y,\mathbb Q_\ell(i+b)),
\end{equation}
basta aplicar el Lema~1.1 a la situación
\begin{equation}\tag{1.2.11}
\begin{tikzcd}[column sep=large,row sep=large]
X\times_S Y \arrow[rr,"\mathrm{pr}_2"]
 \arrow[dr,"\sisi{f\,\mathrm{pr}_1}{f\circ\mathrm{pr}_1}"']
 & {} & Y \arrow[dl,"g"] \\
{} & S & {}
\end{tikzcd}
\qquad .
\end{equation}

\subsubsection*{Corolario 1.4.}
En la situación (1.2.1)--(1.2.2), el morfismo (1.2.3)
\begin{equation}\tag{1.4.1}
z:H^q(X,\mathbb Q_\ell(i))\longrightarrow
H^{q+a-2d_1}(Y,\mathbb Q_\ell(i+b-d_1))
\end{equation}
transforma clases primitivas (XVIII~5.8.1) en clases primitivas, y el
diagrama siguiente es conmutativo:

\pageoriginale

\begin{equation}\tag{1.4.2}
\begin{tikzcd}[column sep=large,row sep=large]
\mathrm{Prim}^q(X/S,\mathbb Q_\ell(i))
 \arrow[r,"z"] \arrow[d,"u_f"]
& \mathrm{Prim}^{q+a-2d_1}(Y/S,\mathbb Q_\ell(i+b-d_1))
 \arrow[d,"u_g"] \\
H^1(S,R^{q-1}f_*\mathbb Q_\ell(i)) \arrow[r,"z"]
& H^1(S,R^{q+a-2d_1-1}g_*\mathbb Q_\ell(i+b-d_1)).
\end{tikzcd}
\end{equation}
Aquí $u_f$ y $u_g$ son los homomorfismos canónicos de XVIII~5.8.1.

\subsubsection*{1.5.}
Sea $\bar\eta$ un punto geométrico de $S$ (sobre un punto $\eta$ de
$S$), y sea
\begin{equation}\tag{1.5.1}
z_{\bar\eta}\in
H^a\bigl(X_{\bar\eta}\times_{k(\bar\eta)}Y_{\bar\eta},
\mathbb Q_\ell(b)\bigr)
\end{equation}
la clase inducida por $z$ sobre la fibra geométrica correspondiente de
\begin{equation}\tag{1.5.2}
X\times_S Y\longrightarrow S.
\end{equation}

\subsubsection*{Proposición 1.5.3.}
Supongamos que $f$ y $g$ son propios, que $S$ es
\sisi{anexo}{conexo}, y sea $q\in\mathbb Z$. Si la clase $z_{\bar\eta}$
induce el homomorfismo nulo
\begin{equation}\tag{1.5.3}
0=z_{\bar\eta}:H^{q-1}(X_{\bar\eta},\mathbb Q_\ell(i))
\longrightarrow
H^{q+a-2d_1-1}(Y_{\bar\eta},\mathbb Q_\ell(i+b-d_1)),
\end{equation}
entonces, para todo
$\sisi{X}{x}\in\mathrm{Prim}^q(X/S,\mathbb Q_\ell(i))$, se tiene
\begin{equation}\tag{1.5.4}
u_g(z(x))=0\quad\text{en}\quad
H^1\bigl(S,R^{q+a-2d_1-1}g_*\mathbb Q_\ell(i+b-d_1)\bigr).
\end{equation}

\pageoriginale

\noindent\underline{Demostración.}
Por la conmutatividad de (1.4.2), basta ver que $z$ induce el
homomorfismo nulo
\begin{equation}\tag{1.5.5}
z:R^{q-1}f_*\mathbb Q_\ell(i)
\stackrel{0}{\longrightarrow}
R^{q+a-2d_1-1}g_*\mathbb Q_\ell(i+b-d_1).
\end{equation}
En efecto, como $f$ y $g$ son propios y lisos, los haces
$R^\bullet f_*\mathbb Q_\ell(\mathord\cdot)$ y
$R^\bullet g_*\mathbb Q_\ell(\mathord\cdot)$ son constantes torcidos
(SGA~4, XVI~2.2), luego equivalen a sus valores en $\bar\eta$, vistos
como módulos bajo $\pi_1(S,\bar\eta)$.

\subsubsection*{1.6. Una construcción de Manin}

Supongamos dadas dos secciones de $f:X\longrightarrow S$:
\begin{equation}\tag{1.6.0}
\begin{tikzcd}[column sep=huge]
S \arrow[r,bend left=45,"t_1"]
  \arrow[r,bend right=45,"t_2"']
& X \arrow[l,"f" description]
\end{tikzcd}
\qquad ,
\end{equation}
y asociemos a cada una su clase de cohomología
\begin{equation}\tag{1.6.1}
\mathrm{cl}(t_i)=t_{i*}(1)
\in H^{2d_1}(X,\mathbb Q_\ell(d_1)).
\end{equation}
Si las fibras de $f:X\longrightarrow S$ son geométricamente conexas
(más generalmente, si para $s\in S$ los puntos
$\sisi{f(s)\text{ y }g(s)}{t_1(s)\text{ y }t_2(s)}$ pertenecen a la
misma componente conexa de $X_s$), es claro que
\begin{equation}\tag{1.6.2}
\mathrm{cl}(t_1)-\mathrm{cl}(t_2)
\in\mathrm{Prim}^{2d_1}(X/S,\mathbb Q_\ell(d_1)).
\end{equation}

\pageoriginale

Denotamos por
\begin{equation}\tag{1.6.3}
u_f(t_1-t_2)\in H^1(S,R^{2d_1-1}f_*\mathbb Q_\ell(d_1))
\end{equation}
su imagen por
\begin{equation}\tag{1.6.4}
u_f:\mathrm{Prim}^{2d_1}(X/S,\mathbb Q_\ell(d_1))
\longrightarrow H^1(S,R^{2d_1-1}f_*\mathbb Q_\ell(d_1)).
\end{equation}
(Manin estudia en [5] un análogo $\mu(t_1-t_2)$ de este elemento.)

\subsubsection*{1.6.5.}
Sea ahora $z$ un ciclo algebraico sobre $X\times_S Y$, de codimensión al
menos $b$ sobre cada fibra de $X\times_S Y$, que determina una clase de
cohomología (SGA~5, Exposé~IV; la referencia ulterior queda en blanco en
la impresión):
\begin{equation}\tag{1.6.6}
\mathrm{cl}(z)\in H^{2b}(X\times_S Y,\mathbb Q_\ell(b)).
\end{equation}
Esta clase define una aplicación (cf.~(1.2.3))
\begin{equation}\tag{1.6.7}
\mathrm{cl}(z)\mathbin{\cdot}H^{2d_1}(X,\mathbb Q_\ell(d_1))
\longrightarrow H^{2b}(Y,\mathbb Q_\ell(b)),
\end{equation}
y, si están definidos los ciclos $z(t_1)$ y $z(t_2)$, la compatibilidad
del producto cup con las intersecciones de ciclos (SGA~5, Exposé~IV; la
referencia ulterior queda en blanco en la impresión) da
\footnote{Si $S$ es liso sobre un esquema $T$ (por ejemplo,
$T=\operatorname{Spec}k$, con $k$ un cuerpo), bastaría suponer que $z$
tiene codimensión al menos $b$ sobre cada fibra por encima de $T$.}

\pageoriginale

\begin{equation}\tag{1.6.8}
\mathrm{cl}(z)\bigl(\mathrm{cl}(t_1)-\mathrm{cl}(t_2)\bigr)
=\mathrm{cl}\bigl(z(t_1)-z(t_2)\bigr).
\end{equation}
La conmutatividad de (1.4.2) nos da
\begin{equation}\tag{1.6.9}
\mathrm{cl}(z)\bigl(u_f(t_1-t_2)\bigr)
=u_g\bigl(\mathrm{cl}(z(t_1)-z(t_2))\bigr).
\end{equation}
Como corolario de la Proposición~1.5.3 obtenemos:

\subsubsection*{Proposición 1.7.}
Sean $t_1$ y $t_2$ dos secciones de $f:X\longrightarrow S$, suponiendo
que $f$ tiene fibras conexas, y sea $z$ un ciclo algebraico sobre
$X\times_S Y$ de codimensión $b$, tal que esté definido el ciclo
algebraico
\begin{equation}\tag{1.7.1}
z(t_1)-z(t_2).
\end{equation}
Entonces su clase de cohomología es primitiva (XVIII~5.8.1). Si el ciclo
algebraico $\sisi{Z_{\bar\eta}}{z_{\bar\eta}}$ sobre
$X_{\bar\eta}\times Y_{\bar\eta}$ induce el homomorfismo nulo
\begin{equation}\tag{1.7.2}
z_{\bar\eta}:H^{2d_1-1}(X_{\bar\eta},\mathbb Q_\ell(d_1))
\stackrel{0}{\longrightarrow}
H^{2b-1}(Y_{\bar\eta},\mathbb Q_\ell(b)),
\end{equation}
entonces
\begin{equation}\tag{1.7.3}
u_g\bigl(z(t_1)-z(t_2)\bigr)=0\quad\text{en}\quad
H^1(S,R^{2b-1}g_*\mathbb Q_\ell(b)),
\end{equation}
donde $u_g$ es el homomorfismo de XVIII~5.8.1 que interviene en 1.4.

\pageoriginale

\section*{2. Recordatorio sobre el nivel}
\noindent(cf.~[2], 9.7 y 10.1).

\subsubsection*{2.0.1.}
Sea $X$ propio y liso sobre un cuerpo algebraicamente cerrado $k$. Vimos
en (1.2.3) que, para todo
\begin{equation}\tag{2.0.2}
T\quad\text{propio, liso y conexo sobre $k$, de dimensión $d$},
\end{equation}
y todo
\begin{equation}\tag{2.0.3}
\text{ciclo algebraico $z$ de codimensión $b$ sobre $T\times X$},
\end{equation}
la clase de cohomología
\begin{equation}\tag{2.0.4}
\mathrm{cl}(z)\in H^{2b}(T\times X,\mathbb Q_\ell(b))
\end{equation}
da lugar a las aplicaciones
\begin{equation}\tag{2.0.5}
\mathrm{cl}(z):H^q(T,\mathbb Q_\ell(i))\longrightarrow
H^{q+2b-2d}(X,\mathbb Q_\ell(i+b-d)).
\end{equation}
La \underline{filtración por conivel} $(N^j)$ de
$H^q(X,\mathbb Q_\ell(i))$ se define por
\begin{equation}\tag{2.0.6}
N^jH^q(X,\mathbb Q_\ell(i))
=\sum_{(T,z)}\operatorname{Im}\!\left(
\mathrm{cl}(z):H^{q-2j}(T,\mathbb Q_\ell(i-j))
\longrightarrow H^q(X,\mathbb Q_\ell(i))\right),
\end{equation}
donde la suma se extiende a los pares
\begin{equation}\tag{2.0.7}
(T,z),\qquad
\begin{array}{l}
T/k\text{ propio, liso y conexo, de dimensión }d,\\
z\text{ ciclo algebraico sobre }T\times X\text{ de codimensión }j+d.
\end{array}
\end{equation}
Conviene observar que a veces se habla de la filtración creciente $F^j$
«por el nivel», definida por
\begin{equation}\tag{2.0.8}
F^jH^q(X,\mathbb Q_\ell(i))
=N^{\left[\frac{q-j}{2}\right]}H^q(X,\mathbb Q_\ell(i)).
\end{equation}

\pageoriginale

Reemplazando $(T,z)$ por
$(\mathbb P^1\times T,\{\mathrm{pt}\}\times z)$, se ve que
\begin{equation}\tag{2.0.9}
N^jH^q(X,\mathbb Q_\ell(i))
\subset N^{j-1}H^q(X,\mathbb Q_\ell(i)).
\end{equation}
Evidentemente se tiene:

\subsubsection*{Proposición 2.1.}
Sean $X$ propio, liso y geométricamente conexo sobre un cuerpo $k$,
$\bar k$ una clausura algebraica de $k$, y
$X_{\bar k}=X\otimes_k\bar k$. Entonces
$\operatorname{Gal}(\bar k/k)$, actuando sobre
$H^q(X_{\bar k},\mathbb Q_\ell(i))$, preserva la filtración por conivel.

Argumentos estándar (cf.~[2]) demuestran los enunciados siguientes.

\subsubsection*{Proposición 2.2.}
La filtración por el nivel es invariante por extensión $k'/k$ de
cuerpos algebraicamente cerrados.

\subsubsection*{Proposición 2.3.}
Sean $S$ un esquema reducido e irreducible, de punto genérico $\eta$, y
\begin{equation*}
\pi:X\longrightarrow S
\end{equation*}
un morfismo propio y liso, de fibras geométricamente conexas. Si, para
enteros dados $j$ y $q$, se tiene en la fibra genérica geométrica
$X_{\bar\eta}$
\begin{equation}\tag{2.3.1}
H^q(X_{\bar\eta},\mathbb Q_\ell)
=N^jH^q(X_{\bar\eta},\mathbb Q_\ell),
\end{equation}
entonces existe un abierto no vacío $U\subset S$ tal que, en todo punto
$u\in U$, se tiene en la fibra geométrica $X_{\bar u}$
\begin{equation}\tag{2.3.2}
H^q(X_{\bar u},\mathbb Q_\ell)
=N^jH^q(X_{\bar u},\mathbb Q_\ell).
\end{equation}

\pageoriginale

\section*{3. Aplicación a los pinceles de Lefschetz}

\subsubsection*{3.1.}
A partir de ahora fijamos
\begin{equation}\tag{3.1.0}
k,\quad\text{un cuerpo algebraicamente cerrado},
\end{equation}
y
\begin{equation}\tag{3.1.1}
\begin{aligned}
X\subset\mathbb P^r,\quad
&\text{un $k$-esquema proyectivo, liso y conexo}\\
&\text{de dimensión par $n$, que verifica (LV) (XVIII~5.2.2).}
\end{aligned}
\end{equation}
Tomemos
\begin{equation}\tag{3.1.2}
D=\{H_t\}_{t\in\mathbb P^1},
\quad\text{un pincel de Lefschetz de hipersuperficies de grado $d$},
\end{equation}
de fibra genérica
\begin{equation}\tag{3.1.3}
X_\eta,\qquad \eta=\text{el punto genérico de }\mathbb P^1.
\end{equation}
Pongamos (cf.~XVIII~5.4.4)
\begin{equation}\tag{3.1.4}
m:X_{\bar\eta}\hookrightarrow X\times_k k(\bar\eta),
\quad\text{la inclusión},
\end{equation}
y
\begin{equation}\tag{3.1.5}
E^{n-1}(X_{\bar\eta},\mathbb Q_\ell(i))
=\bigl(m^*H^{n-1}(X,\mathbb Q_\ell(i))\bigr)^\perp
\subset H^{n-1}(X_{\bar\eta},\mathbb Q_\ell(i)).
\end{equation}
Vimos en XVIII~6.7.2 que la \underline{fórmula de Picard--Lefschetz}
implica
\begin{equation}\tag{3.1.6}
\begin{gathered}
\text{si $n$ es par, todo subgrupo de índice finito de }
\operatorname{Gal}(k(\bar\eta)/k(\eta))\\
\text{actúa irreduciblemente sobre }
E^{\sisi{2n-1}{n-1}}(X_{\bar\eta},\mathbb Q_\ell(i)).
\end{gathered}
\end{equation}

\subsubsection*{Proposición 3.2.}
Sea $D=\{H_t\}$ un pincel de Lefschetz, y supongamos que $n=2m$ es par.
Entonces, o bien

\pageoriginale

\begin{equation}\tag{3.2.1}
H^{n-1}(X_{\bar\eta},\mathbb Q_\ell(i))
\subset m^*H^{n-1}(X,\mathbb Q_\ell(i))
+N^{m-1}H^{n-1}(X_{\bar\eta},\mathbb Q_\ell(i)),
\end{equation}
o bien:

\subsubsection*{(3.2.2)}
Para toda curva $C$ propia y lisa sobre $k(\bar\eta)$ y todo ciclo
algebraico $z$ sobre $C\times X_{\bar\eta}$ de codimensión $m$, el
homomorfismo inducido (1.2.3)
\begin{equation}\tag{3.2.3}
\mathrm{cl}(z):H^1(C,\mathbb Q_\ell(1+i-m))
\longrightarrow H^{2m-1}(X_{\bar\eta},\mathbb Q_\ell(i))
\end{equation}
tiene su imagen en $m^*H^{2m-1}(X,\mathbb Q_\ell(i))$.

\noindent\underline{Demostración.}
Supongamos que la imagen de tal homomorfismo (3.2.3) no está contenida
en $m^*H^{n-1}(X,\mathbb Q_\ell(i))$, o, lo que es equivalente, que su
imagen en $E^{n-1}(X_{\bar\eta},\mathbb Q_\ell(i))$ no es nula. Esta
imagen es evidentemente estable bajo la acción de un subgrupo de índice
finito de $\operatorname{Gal}(k(\bar\eta)/k(\eta))$. La irreducibilidad
(3.1.6) de $E^{n-1}(X_{\bar\eta},\mathbb Q_\ell(i))$ dice entonces que
\begin{equation*}
H^{n-1}(X_{\bar\eta},\mathbb Q_\ell(i))
=m^*H^{n-1}(X,\mathbb Q_\ell(i))+\operatorname{Im}(3.2.3).
\end{equation*}
Por definición,
\begin{equation*}
\operatorname{Im}(3.2.3)
=\mathrm{cl}(z)\bigl(H^1(C,\mathbb Q_\ell(1+i-m))\bigr)
\subset N^{m-1}H^{n-1}(X_{\bar\eta},\mathbb Q_\ell(i)).
\end{equation*}
Por consiguiente,
\begin{equation*}
H^{n-1}(X_{\bar\eta},\mathbb Q_\ell(i))
=m^*H^{n-1}(X,\mathbb Q_\ell(i))
+N^{m-1}H^{n-1}(X_{\bar\eta},\mathbb Q_\ell(i)),
\end{equation*}
que es precisamente la primera alternativa.

\pageoriginale

\subsubsection*{3.3. La condición (B)}

\paragraph{3.3.1.}
Diremos que la inclusión
\begin{equation*}
m:X_{\bar\eta}\hookrightarrow X\times_k k(\bar\eta)
\end{equation*}
verifica la condición (B) si existe
\begin{equation}\tag{3.3.2}
\begin{gathered}
V_{\bar\eta},\quad\text{un ciclo algebraico de codimensión $n-1$ sobre}\\
X_{\bar\eta}\times\bigl(X\times_k k(\bar\eta)\bigr),
\end{gathered}
\end{equation}
tal que la aplicación que define
\begin{equation}\tag{3.3.3}
V_{\bar\eta}:H^{n-1}(X_{\bar\eta},\mathbb Q_\ell(i))
\longrightarrow H^{n-1}(X,\mathbb Q_\ell(i))
\end{equation}
sea inversa por la izquierda de
\begin{equation}\tag{3.3.4}
m^*:H^{n-1}(X,\mathbb Q_\ell(i))
\longrightarrow H^{n-1}(X_{\bar\eta},\mathbb Q_\ell(i)),
\end{equation}
es decir, tal que
\begin{equation}\tag{3.3.5}
V_{\bar\eta}\,m^*=\mathrm{id}.
\end{equation}

\subsubsection*{3.4.}
La condición (B) se verifica si la operación $\wedge$ de [4, 1.4.2.1]
es algebraica (cf.~[4, 2.12(ii)]). Se conjetura que $\wedge$ es siempre
algebraica y, por tanto, que la condición (B) se verifica siempre; esto
se sabe comprobar en numerosos casos [4]. Se verifica trivialmente si
$X$ es una intersección completa de dimensión par, pues entonces se sabe
(IX~1.6) que
\begin{equation}\tag{3.4.1}
H^{n-1}(X,\mathbb Q_\ell(i))=0.
\end{equation}

\pageoriginale

\section*{4. El teorema de Griffiths}

\subsubsection*{Teorema 4.1.}
Sea $X$ proyectivo, liso y conexo sobre $k$, de dimensión par
$n=2m\geq4$. Sea $D$ un pincel de Lefschetz de secciones por
hipersuperficies de grado $d$. Supongamos que:

\medskip
\noindent\textbf{(4.1.0)} $X$ verifica (LV) (XVIII~5.2.2);

\noindent\textbf{(4.1.1)} el pincel $D$ verifica la hipótesis (A)
(XVIII~5.3);

\noindent\textbf{(4.1.2)} el pincel $D$ verifica la alternativa (3.2.2),
es decir, en la situación de la Proposición~3.2, la cohomología
evanescente $E^{n-1}(X_{\bar\eta},\mathbb Q_\ell)$ no es de nivel
$\leq1$;

\noindent\textbf{(4.1.3)} la inclusión
$\sisi{m^*}{m}:X_{\bar\eta}\hookrightarrow X\times_k k(\bar\eta)$
verifica la hipótesis (B) (3.3).

\medskip
Entonces, si un ciclo algebraico $z$ de codimensión $m$ sobre $X$ tiene
una restricción a $X_{\bar\eta}$ que, como ciclo algebraico, satisface
\begin{equation}\tag{4.1.4}
z_{\bar\eta}=z|_{X_{\bar\eta}}
\quad\text{es algebraicamente equivalente a cero},
\end{equation}
entonces
\begin{equation}\tag{4.1.5}
\operatorname{cl}(z)=0
\quad\text{en }H^{2m}(X,\mathbb Q_\ell(m));
\end{equation}
en otras palabras, $z$ es $\mathbb Q_\ell$-homólogo a cero.

\subsubsection*{\texorpdfstring{4.2. Una observación sobre el caso en que $X$ es una
intersección completa}{4.2. Una observación sobre el caso en que X es una intersección completa}}

Como ya se ha observado (cf.~XVIII~5.2.5 y (3.4)), las hipótesis (4.1.0)
y (4.1.3) se verifican entonces automáticamente. Vimos en XVIII~6.3.4 y
6.4 que la hipótesis (4.1.1) se verifica para un pincel de
hipersuperficies de grado $d>2n$. Por otra parte, en la exposición siguiente
probaremos la condición (4.1.2) para un pincel «general» de
hipersuperficies de grado $d>2n$ (XXI~5.2). Por consiguiente, un pincel
«general» de hipersuperficies de grado $d>2n$ verifica todas las
hipótesis del Teorema~4.1.

\pageoriginale

\subsubsection*{4.3.}
Volviendo al caso de $X$ arbitrario, parece muy plausible que, para un
grado $d$ suficientemente grande, todas las condiciones de 4.1 se
verifiquen para $D$, al menos si $D$ es suficientemente general. Ya lo
hemos visto para (4.1.1) (XVIII~6.4), y hemos señalado la conjetura de que
(4.1.0) y (4.1.3) se verifican siempre (siendo cierto en todo caso para
(4.1.0) cuando $\operatorname{char}(k)=0$). En cuanto a (4.1.2), la teoría
de Hodge la demuestra para $d$ grande si
$\operatorname{char}(k)\sisi{\ne}{=}0$, mientras que argumentos
heurísticos, que el lector encontrará en el apéndice a la exposición siguiente,
hacen plausible que ocurra lo mismo si $\operatorname{char}(k)>0$.

Más generalmente, estos argumentos ---la teoría de Hodge y,
respectivamente, la cohomología cristalina--- demuestran, y
respectivamente hacen plausible, que para una variedad proyectiva lisa
$X$ de dimensión $n$, y para todo pincel de Lefschetz $D$ de
hipersuperficies de grado $d$ suficientemente grande (respectivamente,
todo pincel $D$ suficientemente general de hipersuperficies de grado
$d$), la cohomología evanescente
$E^{n-1}(X_{\bar\eta},\mathbb Q_\ell)$ es de nivel exactamente $n-1$;
es decir, no es de nivel $\leq n-3$, o equivalentemente no es de conivel
$1$.

\subsubsection*{4.4. Demostración del teorema de Griffiths}

Sea $z$ un ciclo algebraico de codimensión $\sisi{n}{m}$ sobre $X$ tal
que
\begin{equation}\tag{4.4.1}
z_{\bar\eta}=z|_{X_{\bar\eta}}
\quad\text{es algebraicamente equivalente a cero}.
\end{equation}
Por definición, existen
\begin{equation}\tag{4.4.2}
\text{una curva $C_{\bar\eta}$ propia, lisa y conexa sobre $k(\bar\eta)$},
\end{equation}

\pageoriginale

\begin{equation}\tag{4.4.3}
t_i(\bar\eta):\operatorname{Spec}k(\bar\eta)
\longrightarrow\sisi{C}{C_{\bar\eta}},\qquad i=1,2,
\end{equation}
y
\begin{equation}\tag{4.4.4}
\text{un ciclo algebraico $P_{\bar\eta}$ de codimensión
$\sisi{n}{m}$ sobre $C_{\bar\eta}\times X_{\bar\eta}$},
\end{equation}
tales que
\begin{equation}\tag{4.4.5}
z_{\bar\eta}
=P_{\bar\eta}\bigl(t_1(\bar\eta)\bigr)
-P_{\bar\eta}\bigl(t_2(\bar\eta)\bigr)
\quad\text{sobre }X_{\bar\eta}.
\end{equation}
Consideremos $P_{\bar\eta}$ como una correspondencia algebraica de
$C_{\bar\eta}$ hacia $X_{\bar\eta}$. Define una correspondencia
cohomológica, denotada aún por $P_{\bar\eta}$, o equivalentemente un
homomorfismo
\begin{equation}\tag{4.4.6}
P_{\bar\eta}:H^*(C_{\bar\eta},\mathbb Q_\ell(i))
\longrightarrow
H^{*+2(m-1)}\bigl(X_{\bar\eta},
 \mathbb Q_\ell(i+m-1)\sisi{}{\bigr)}.
\end{equation}
Consideremos asimismo la correspondencia algebraica $V_{\bar\eta}$ de
(3.3.2), de $X_{\bar\eta}$ hacia $X\otimes_k k(\bar\eta)$, dada por la
hipótesis (B) de (4.1.3). Define
\begin{equation}\tag{4.4.7}
V_{\bar\eta}:H^*(X_{\bar\eta},\mathbb Q_\ell(i+m-1))
\longrightarrow
H^*(X\otimes_k k(\bar\eta),\mathbb Q_\ell(i+m-1)).
\end{equation}
La correspondencia algebraica compuesta
$Q_{\bar\eta}=V_{\bar\eta}\circ P_{\bar\eta}$, de $C_{\bar\eta}$ hacia
$X\otimes_k k(\bar\eta)$, define por tanto
\begin{equation}\tag{4.4.8}
Q_{\bar\eta}=V_{\bar\eta}\circ P_{\bar\eta}:
H^*(C_{\bar\eta},\mathbb Q_\ell(i))
\longrightarrow
H^{*+2(m-1)}(X\otimes_k k(\bar\eta),\mathbb Q_\ell(i+m-1)).
\end{equation}
Consideremos la inclusión
\begin{equation*}
m:X_{\bar\eta}\longrightarrow X\otimes_k k(\bar\eta).
\end{equation*}
Da la correspondencia algebraica $m^*$ de
$X\otimes_k k(\bar\eta)$ hacia $X_{\bar\eta}$, y por tanto la
correspondencia compuesta
$m^*Q_{\bar\eta}\sisi{\,{}_{\bar\eta}}{}$ de $C_{\bar\eta}$ hacia
$X_{\bar\eta}$, que define

\pageoriginale

\begin{equation}\tag{4.4.9}
m^*Q_{\bar\eta}=m^*V_{\bar\eta}P_{\bar\eta}:
H^*(C_{\bar\eta},\mathbb Q_\ell(i))
\longrightarrow
H^{*+2(m-1)}(X_{\bar\eta},\mathbb Q_\ell(i+m-1)).
\end{equation}
Consideremos ahora la correspondencia algebraica
\begin{equation}\tag{4.4.10}
\begin{split}
R_{\bar\eta}
&=P_{\bar\eta}-m^*Q_{\bar\eta}
=(\mathrm{id}-m^*V_{\bar\eta})P_{\bar\eta}:\\
&\hspace{2em}H^*(C_{\bar\eta},\mathbb Q_\ell(i))
\longrightarrow
H^{*+2(m-1)}(X_{\bar\eta},\mathbb Q_\ell(i+m-1)).
\end{split}
\end{equation}
Como $V_{\bar\eta}$ es inversa por la izquierda de $m^*$,
$\mathrm{id}-m^*V_{\bar\eta}$ es un proyector sobre
$H^*(X_{\bar\eta},-)$ cuya imagen es un suplementario de
\begin{equation*}
\operatorname{Im}\bigl(m^*:H^*(X\otimes_k k(\bar\eta),-)
\longrightarrow H^*(X_{\bar\eta},-)\bigr).
\end{equation*}
Por la hipótesis (4.1.2), la imagen de
\begin{equation}\tag{4.4.11}
R_{\bar\eta}:H^1(C_{\bar\eta},\mathbb Q_\ell(i))
\longrightarrow
H^{2m-1}(X_{\bar\eta},\mathbb Q_\ell(i+m-1))
\end{equation}
está contenida en $\operatorname{Im}m^*$. Por consiguiente,
\begin{equation}\tag{4.4.12}
\text{el homomorfismo (4.4.11) es nulo.}
\end{equation}
Elijamos una extensión finita
\begin{equation*}
L/k(\eta)
\end{equation*}
sobre la cual los objetos
\begin{equation*}
C_{\bar\eta},\quad t_1(\bar\eta),\quad t_2(\bar\eta),\quad
Q_{\bar\eta},\quad R_{\bar\eta}
\end{equation*}
estén ya definidos, y sea $U$ un abierto no vacío de la normalización de
$\mathbb P^1$ en $L$. Encogiendo $U$, existe una curva relativa propia y
lisa, de fibras geométricamente conexas,
\begin{equation*}
C\longrightarrow U,
\end{equation*}
cuya fibra genérica geométrica sea $C_{\bar\eta}$, junto con dos
secciones
\begin{equation*}
t_i:U\longrightarrow C,\qquad i=1,2,
\end{equation*}
cuyos valores en $\bar\eta$ son $t_i(\bar\eta)$.

\pageoriginale

Sea
\begin{equation*}
\rho:\widetilde X\longrightarrow\mathbb P^1
\end{equation*}
la proyección de XVIII~3.1.4 cuyas fibras son los miembros del pincel,
pongamos
\begin{equation*}
\widetilde X_U=\widetilde X\times_{\mathbb P^1}U,
\end{equation*}
y sea
\begin{equation*}
M:\widetilde X_U\hookrightarrow X\times U
\end{equation*}
la inmersión canónica en la familia constante sobre $U$ de valor $X$.
Pongamos
\begin{equation*}
z_U=M^*\mathrm{pr}_1^*(z),
\end{equation*}
el ciclo algebraico sobre $\widetilde X_U$ inducido por $z$.

Encogiendo de nuevo $U$, existen un ciclo algebraico $R$ sobre
$C\times_U\widetilde X_U$, cuyo valor sobre
$C_{\bar\eta}\times X_{\bar\eta}$ es $R_{\bar\eta}$, y un ciclo
algebraico $Q$ sobre $C\times X\times U$, cuyo valor sobre
$C_{\bar\eta}\times(X\otimes_k k(\bar\eta))$ es $Q_{\bar\eta}$, tales
que
\begin{equation}\tag{4.4.13}
z_U=R(t_1)-R(t_2)
+(\mathrm{id}\times M)^*\bigl[Q(t_1)-Q(t_2)\bigr].
\end{equation}
Para demostrar el teorema basta, gracias a la inyectividad del
homomorfismo de Griffiths «módulo la parte fija» (XVIII~5.8.7),

\pageoriginale

\begin{equation*}
\mathrm{Prim}^{2m}\bigl(X,\mathbb Q_\ell(\sisi{n}{m})\bigr)
\xrightarrow{\ \underline{\operatorname{grif}}\ }
\frac{H^1(U,R^{2m-1}\rho_{U*}\mathbb Q_\ell(m))}
{M^*H^1(U,H^{2m-1}(X,\mathbb Q_\ell(m)))},
\end{equation*}
probar que la imagen de $z_U$ por
\begin{equation*}
u_0:\mathrm{Prim}^{2m}(\widetilde X_U/U,\mathbb Q_\ell(m))
\longrightarrow H^1(U,R^{2m-1}\rho_{U*}\mathbb Q_\ell(m))
\end{equation*}
pertenece a
$M^*H^1(U,H^{2m-1}(X,\mathbb Q_\ell(m)))$.

En vista de la descomposición (4.4.13) y de la conmutatividad de (1.4.2),
tenemos el diagrama conmutativo
\begin{center}
\small
\begin{tikzcd}[column sep=2.2em,row sep=3em]
\mathrm{Prim}^{2}(C/U,\mathbb Q_\ell(1))
 \arrow[r,"Q"]
& \mathrm{Prim}^{2m}(X\times U/U,\mathbb Q_\ell(\sisi{n}{m}))
 \arrow[r,"M^*"] \arrow[d,"u_{\mathrm{pr}_2}"]
& \mathrm{Prim}^{2m}(\widetilde X_U/U,\mathbb Q_\ell(m))
 \arrow[d,"u_0"] \\
{}
& H^1(U,H^{2m-1}(X,\mathbb Q_\ell(m))\sisi{}{)}
 \arrow[r,"M^*"]
& H^1(U,R^{2m-1}\rho_{U*}\mathbb Q_\ell(m)),
\end{tikzcd}
\end{center}
por lo que basta demostrar
\begin{equation*}
u_0\bigl(R(t_1)-R(t_2)\bigr)=0
\quad\text{en}\quad
H^1\bigl(U,R^{\sisi{2n-1}{2m-1}}\rho_{U*}
 \mathbb Q_\ell(\sisi{n}{m})\bigr).
\end{equation*}
Esto resulta de $\sisi{(4.3.12)}{(4.4.12)}$, por 1.7.

\section*{5. El grupo de Griffiths}

\subsubsection*{5.0.}
Sea $S$ propio y liso sobre un cuerpo algebraicamente cerrado $k$.
Denotemos por

\pageoriginale

\medskip
\noindent\textbf{(5.0.1)} $Z^i(S)$ el grupo de los ciclos algebraicos de
codimensión $i$ sobre $S$;

\noindent\textbf{(5.0.2)} $Z^i_{\mathrm{alg}}(S)$ el subgrupo de $Z^i(S)$
formado por los ciclos algebraicamente equivalentes a cero; y, para todo
número primo $\ell\ne\operatorname{char}(k)$,

\noindent\textbf{(5.0.3)} $Z^i_{\ell\text{-}\mathrm{coh}}(S)$ el
subgrupo de $Z^i(S)$ formado por los ciclos cuya clase de cohomología en
$H^{2i}(S,\mathbb Q_\ell(i))$ es nula.

Se conjetura en [4] que este grupo es independiente de $\ell$; esto es
cierto al menos si $\operatorname{char}(k)=0$. Pongamos
\begin{equation}\tag{5.0.4}
\mathrm{Grif}^i_\ell(S)
=Z^i_{\ell\text{-}\mathrm{coh}}(S)/Z^i_{\mathrm{alg}}(S).
\end{equation}
Si además $S$ es proyectivo y está provisto de una inmersión proyectiva,
pongamos

\medskip
\noindent\textbf{(5.0.5)} $Z^{i,\mathrm{prim}}(S)$ igual al subgrupo de
$Z^i(S)$ formado por los ciclos primitivos (es decir, aquellos cuya
clase de cohomología es primitiva), y

\noindent\textbf{(5.0.6)}
$\mathrm{Prim}^{2i}_{\mathrm{alg}}(S,\mathbb Q_\ell(i))$ igual al
subespacio vectorial sobre $\mathbb Q$ de
$H^{2i}(S,\mathbb Q_\ell(i))$ engendrado por las clases de cohomología de
los ciclos algebraicos primitivos.

\subsubsection*{Corolario 5.1.}
Bajo las hipótesis del Teorema~4.1, se tiene
\begin{equation}\tag{5.1.1}
\dim_{\mathbb Q}\!\left(
 \mathrm{Grif}^{\sisi{n}{m}}_\ell(X_{\bar\eta})
 \otimes_{\mathbb Z}\mathbb Q\right)
\geq
\dim_{\mathbb Q}\!\left(
 \mathrm{Prim}^{\sisi{2n}{2m}}_{\mathrm{alg}}
 (X,\mathbb Q_\ell(\sisi{n}{m}))\right).
\end{equation}
En efecto, por (4.1.5), el núcleo de la aplicación compuesta
\begin{equation*}
\mathbb Q\otimes_{\mathbb Z}Z^{\sisi{n}{m},\mathrm{prim}}(X)
\xrightarrow{\ \text{restr. a }X_{\bar\eta}\ }
\mathbb Q\otimes_{\mathbb Z}
 Z^{\sisi{n}{m}}_{\ell\text{-}\mathrm{coh}}(X_{\bar\eta})
\longrightarrow
\mathrm{Grif}^{\sisi{n}{m}}_\ell(X_{\bar\eta})
\end{equation*}
está contenido en
$Z^{\sisi{n}{m}}_{\ell\text{-}\mathrm{coh}}(X)
\otimes_{\mathbb Z}\mathbb Q$. Por consiguiente,
$\mathrm{Prim}^{\sisi{2n}{2m}}_{\mathrm{alg}}
(X,\mathbb Q_\ell(\sisi{n}{m}))$ es un cociente de un subespacio de
$\mathrm{Grif}^{\sisi{n}{m}}_\ell(X_{\bar\eta})
\otimes_{\mathbb Z}\mathbb Q$, lo que prueba la afirmación.

\pageoriginale

\subsubsection*{Observación 5.2.}
Para $X$ una hipersuperficie cuádrica en $\mathbb P^{2m+1}$ se tiene
\begin{equation*}
\dim_{\mathbb Q_\ell}\!\bigl(
\mathrm{Prim}^{2m}_{\mathrm{alg}}
(X,\mathbb Q_\ell(\sisi{n}{m}))\sisi{}{\bigr)}=1.
\end{equation*}
En particular, $\sisi{\text{(4.2), }}{}$cuando $m=2$ y $X$ es una
cuádrica en $\mathbb P^5$, la sección de $X$ por una hipersuperficie
genérica de grado $>5$ contiene una curva $\mathbb Q_\ell$-homóloga a
cero de la cual ningún múltiplo no nulo es algebraicamente equivalente
a cero.

\begin{thebibliography}{9}
\bibitem{GriffithsPeriodsXX}
P. Griffiths, «On the Periods of Certain Rational Integrals III», de
próxima aparición en \emph{Annals of Mathematics}.

\bibitem{GrothendieckBrauerIIIXX}
A. Grothendieck, «Le Groupe de Brauer III» (especialmente
\S\S~9--10), en \emph{Dix Exposés sur la Cohomologie des Schémas},
North-Holland, 1968.

\bibitem{JouanolouCohomologieXX}
J.-P. Jouanolou, \emph{Cohomologie $\ell$-adique}, tesis, París, 1969.

\bibitem{KleimanCyclesXX}
S. Kleiman, «Algebraic Cycles and the Weil Conjectures», en
\emph{Dix Exposés sur la Cohomologie des Schémas}, North-Holland, 1968.

\bibitem{ManinCurvesXX}
J. Manin, «Rational Points of Algebraic Curves over Function Fields»,
\emph{Translations of the American Mathematical Society}.
\end{thebibliography}
